%% Guided notes, AP Calculus AB Unit 1: Limits, Continuity, and Differentiability (lessons 1-1 to 1-9).
\newcommand{\gr}[2]{\begin{center}\begin{tikzpicture}\begin{axis}[gridaxis,#1]#2\end{axis}\end{tikzpicture}\end{center}}
\newcommand{\pt}[1]{\node[circle,fill=apblue,inner sep=1.5pt] at (axis cs:#1){};}
\newcommand{\op}[1]{\node[circle,draw=apblue,fill=white,thick,inner sep=1.5pt] at (axis cs:#1){};}
\newcommand{\va}[1]{\draw[apred,dashed,thick](axis cs:#1,-100)--(axis cs:#1,100);}
\newcommand{\ha}[1]{\draw[apgold!80!black,dashed,thick](axis cs:-100,#1)--(axis cs:100,#1);}
\newcommand{\fgraph}{\gr{xmin=-3.5,xmax=4.5,ymin=-0.5,ymax=4.5,width=7.5cm,height=5cm}{\addplot[apblue,very thick,domain=-3:-1]{x+4};\addplot[apblue,very thick,domain=-1:2]{x^2};\addplot[apblue,very thick,domain=2:4]{8-2*x};\op{-1,3}\pt{-1,1}\op{2,4}\pt{2,1}\pt{-3,1}\pt{4,0}\node[font=\small,apblue] at (axis cs:3.6,3.4){$y=f(x)$};}}

\NewEnviron{nrow}[1]{\begin{tcolorbox}[enhanced,sidebyside,sidebyside align=top seam,lefthand width=2.7cm,sidebyside gap=12pt,colback=white,colframe=apnavy!55,boxrule=0.5pt,arc=0pt,left=5pt,right=6pt,top=5pt,bottom=5pt,before skip=0pt,after skip=-0.5pt,segmentation style={solid,apnavy!35,line width=0.5pt},fontupper=\sffamily\bfseries\color{apnavy}\raggedright]#1\tcblower\BODY\end{tcolorbox}}
\newcommand{\sg}[2]{\begin{center}\begin{tikzpicture}\begin{axis}[gridaxis,#1]#2\end{axis}\end{tikzpicture}\end{center}}
\newcommand{\tinyg}[2]{\begin{tikzpicture}\begin{axis}[gridaxis,width=4.6cm,height=4cm,xticklabels={},yticklabels={},#1]#2\end{axis}\end{tikzpicture}}

%%DOC 1-1 | Limits to a Point on a Graph
\begin{center}\textsf{\small\color{apgrey}Objective: find left-hand, right-hand and two-sided limits from a graph.}\end{center}
\begin{nrow}{Limit notation}
\begin{tabular}{@{}p{0.31\linewidth}p{0.31\linewidth}p{0.3\linewidth}@{}}
\textbf{From the left} & \textbf{From the right} & \textbf{From both sides}\\[2pt]
\fillin{$\displaystyle\lim_{x\to c^-}f(x)$} & \fillin{$\displaystyle\lim_{x\to c^+}f(x)$} & \fillin{$\displaystyle\lim_{x\to c}f(x)$}\\[4pt]
\footnotesize $x$ approaches $c$ with values \fillin{less} than $c$ & \footnotesize $x$ approaches $c$ with values \fillin{greater} than $c$ & \footnotesize both one-sided limits are \fillin{equal}
\end{tabular}
\end{nrow}
\begin{nrow}{My first limits}
\sg{xmin=-1.5,xmax=6.5,ymin=-0.5,ymax=6.5,width=8.5cm,height=5.6cm,xtick={-1,0,1,2,3,5,6},extra x ticks={4},extra x tick labels={$c$}}{\addplot[apblue,very thick,domain=-1:6,samples=80]{0.25*(x-2)^2+2};\draw[apred,dashed](axis cs:4,0)--(axis cs:4,3)--(axis cs:0,3);\pt{4,3}\node[apblue,font=\small] at (axis cs:5.6,4.8){$f(x)$};}
Use the graph, with $c=4$.\par
1) $\displaystyle\lim_{x\to4^-}f(x)=$ \fillin{$3$} \quad 2) $\displaystyle\lim_{x\to4^+}f(x)=$ \fillin{$3$} \quad 3) $\displaystyle\lim_{x\to4}f(x)=$ \fillin{$3$}\par
\smallskip A limit is a \fillin{$y$-value} (a height): the height the graph \fillin{approaches} as $x$ gets close to $c$.
\end{nrow}
\begin{nrow}{Challenge}
Draw the graph of a function whose left-hand and right-hand limits at $x=1$ are \textbf{not} the same.
\sg{xmin=-4.5,xmax=4.5,ymin=-3.5,ymax=4.5,width=7cm,height=5.2cm}{\ifkey\addplot[apblue,very thick,domain=-4:1]{x+2};\addplot[apblue,very thick,domain=1:4]{-x+1};\pt{1,3}\op{1,0}\fi}
\ifkey{\color{apnavy}\small Example: the left piece arrives at height $3$, the right piece starts at height $0$, so $\lim_{x\to1}f(x)$ does not exist.}\fi
\end{nrow}
\begin{nrow}{Example:\\ one-sided limits are not the same}
\sg{xmin=-3.5,xmax=5,ymin=-0.5,ymax=4.2,width=8.5cm,height=5.4cm}{\addplot[apblue,very thick,domain=-3:2]{0.5*x+2};\addplot[apblue,very thick,domain=2:4.5,samples=60]{(x-3)^2};\pt{2,3}\op{2,1}\draw[apgrey,dashed](axis cs:2,0)--(axis cs:2,4);\node[apblue,font=\small] at (axis cs:4.3,3.4){$g(x)$};}
1) $\displaystyle\lim_{x\to2^-}g(x)=$ \fillin{$3$} \qquad 2) $\displaystyle\lim_{x\to2^+}g(x)=$ \fillin{$1$}\par
3) $\displaystyle\lim_{x\to2}g(x)=$ \fillin{DNE} \qquad 4) $g(2)=$ \fillin{$3$} \;{\small(the height \emph{at} $x=2$: the filled dot)}\par
\smallskip If the left and right limits are different, the two-sided limit \fillin{does not exist} (DNE).
\end{nrow}
\begin{nrow}{Summary}
\ifkey{\color{apnavy}A limit is the height the graph approaches. $\lim_{x\to c}f(x)$ exists only when the left and right limits are equal. The value $f(c)$ is the height of the filled dot at $c$; it can be different from the limit, or not exist.}\else\vspace{2.6cm}\fi
\end{nrow}

%%DOC 1-2 | Limits to a Point of a Function
\begin{center}\textsf{\small\color{apgrey}Objective: evaluate limits algebraically, including when plugging in fails.}\end{center}
\begin{nrow}{Evaluating limits}
\begin{tabular}{@{}p{0.31\linewidth}p{0.31\linewidth}p{0.31\linewidth}@{}}
\textbf{Linear} $f(x)=2x-1$ & \textbf{Quadratic} $g(x)=x^2-3$ & \textbf{Constant} $h(x)=-2$\\
\tinyg{xmin=-1.5,xmax=4.5,ymin=-3.5,ymax=6.5}{\addplot[apblue,very thick,domain=-1.2:3.7]{2*x-1};\pt{3,5}} &
\tinyg{xmin=-3.5,xmax=3.5,ymin=-3.5,ymax=6.5}{\addplot[apblue,very thick,domain=-3:3,samples=60]{x^2-3};\pt{-2,1}} &
\tinyg{xmin=-1.5,xmax=6.5,ymin=-3.5,ymax=3.5}{\addplot[apblue,very thick,domain=-1.5:6.5]{-2};\pt{5,-2}}\\
$\displaystyle\lim_{x\to3}f(x)=$ \fillin{$5$} & $\displaystyle\lim_{x\to-2}g(x)=$ \fillin{$1$} & $\displaystyle\lim_{x\to5}h(x)=$ \fillin{$-2$}\\[6pt]
$f(3)=$ \fillin{$5$} & $g(-2)=$ \fillin{$1$} & $h(5)=$ \fillin{$-2$}
\end{tabular}\par
\smallskip For polynomials (and constants), the limit equals the value: just \fillin{plug in} $c$.
\end{nrow}
\begin{nrow}{Limits of piecewise functions}
\begin{minipage}[t]{0.56\linewidth}
$j(x)=\begin{cases}x+3, & x\le1\\ x^2+1, & x>1\end{cases}$\par\smallskip
$x=1$ is the \fillin{suspect point}: use the \fillin{top} piece for the left and the \fillin{bottom} piece for the right.\par\smallskip
$\displaystyle\lim_{x\to1^-}j(x)=$ \fillin{$4$} \quad $\displaystyle\lim_{x\to1^+}j(x)=$ \fillin{$2$}\par
$\displaystyle\lim_{x\to1}j(x)=$ \fillin{DNE} \quad $j(1)=$ \fillin{$4$}
\end{minipage}\hfill
\begin{minipage}[t]{0.4\linewidth}\vspace{0pt}
\tinyg{xmin=-3.5,xmax=2.5,ymin=-0.5,ymax=6,width=5.2cm,height=4.8cm}{\ifkey\addplot[apblue,very thick,domain=-3:1]{x+3};\addplot[apblue,very thick,domain=1:2.2,samples=40]{x^2+1};\pt{1,4}\op{1,2}\fi}
\end{minipage}
\end{nrow}
\begin{nrow}{When plugging in fails}
$\displaystyle\lim_{x\to3}\frac{x^2-9}{x^2-x-6}$\quad Plug in: \fillin{$\dfrac{0}{0}$}, the \fillin{indeterminate form}. It means: do more work!\par\smallskip
\begin{minipage}[t]{0.3\linewidth}Three steps:\\ 1) \fillin{Factor}\\ 2) \fillin{Cancel}\\ 3) \fillin{Plug in}\end{minipage}\hfill
\begin{minipage}[t]{0.66\linewidth}
\begin{workspace}[2.4cm]
$\dfrac{(x-3)(x+3)}{(x-3)(x+2)}=\dfrac{x+3}{x+2}\;\to\;\dfrac{3+3}{3+2}=\dfrac65$
\end{workspace}
\end{minipage}
\end{nrow}
\begin{nrow}{Example 1}
$\displaystyle\lim_{x\to0}\frac{x^2-4x}{3x^2+x}$
\begin{workspace}[2.2cm]
$\frac00$. \;$\dfrac{x(x-4)}{x(3x+1)}=\dfrac{x-4}{3x+1}\to\dfrac{-4}{1}=-4$.
\end{workspace}
\end{nrow}
\begin{nrow}{Example 2}
$\displaystyle\lim_{x\to-3}\frac{x^2+6x+9}{x+3}$
\begin{workspace}[2.2cm]
$\frac00$. \;$\dfrac{(x+3)(x+3)}{x+3}=x+3\to-3+3=0$. \;(A limit of $0$ is a real answer, not DNE.)
\end{workspace}
\end{nrow}
\begin{nrow}{Summary}
\ifkey{\color{apnavy}Try plugging in first. A number means you are done; $\frac{0}{0}$ means factor, cancel, then plug in again. For a piecewise function, compare the left piece and the right piece at the suspect point.}\else\vspace{2.4cm}\fi
\end{nrow}

%%DOC 1-3 | Limits to Infinity and Horizontal Asymptotes
\begin{center}\textsf{\small\color{apgrey}Objective: (1) find vertical asymptotes and holes; (2) use the top/bottom rules for limits to infinity and horizontal asymptotes.}\end{center}
\begin{nrow}{The graph}
\sg{xmin=-5.5,xmax=6.5,ymin=-4.5,ymax=6.5,width=11cm,height=6.8cm}{\addplot[apblue,very thick,domain=-5.5:-1.18,samples=80]{(x+2)/(x+1)};\addplot[apblue,very thick,domain=-0.83:6.5,samples=100]{(x+2)/(x+1)};\va{-1}\ha{1}\op{2,1.3333}}
\centering Graph of $f(x)=\dfrac{x^2-4}{x^2-x-2}$
\end{nrow}
\begin{nrow}{Vertical asymptotes (VA)}
Factor: $f(x)=\dfrac{x^2-4}{x^2-x-2}=$ \fillin{$\dfrac{(x-2)(x+2)}{(x-2)(x+1)}$}\par\smallskip
\textbf{Vertical asymptote}: the \fillin{denominator} is $0$ after cancelling: \fillin{$x=-1$}.\par
\textbf{Hole}: a factor that \fillin{cancels}: \fillin{$x=2$}, at height $\frac{2+2}{2+1}=$ \fillin{$\frac43$}. A hole is a \fillin{removable} discontinuity.\par
Near a VA the limit is \fillin{$\pm\infty$}: \;$\displaystyle\lim_{x\to-1^+}f(x)=$ \fillin{$+\infty$}, \;$\displaystyle\lim_{x\to-1^-}f(x)=$ \fillin{$-\infty$}.
\end{nrow}
\begin{nrow}{Horizontal asymptotes (HA)}
A horizontal asymptote comes from a limit to \fillin{infinity}: $\displaystyle\lim_{x\to\infty}f(x)=$ \fillin{$1$}, so the HA is \fillin{$y=1$}.\par\smallskip
Fun fact: for a rational function with a HA, \fillin{$\displaystyle\lim_{x\to\infty}f(x)=\lim_{x\to-\infty}f(x)$}: the same HA on both ends.
\end{nrow}
\begin{nrow}{Limits to infinity: top/bottom rules}
Compare the \fillin{highest powers} (degrees) of the top and the bottom.\par\smallskip
1) $\displaystyle\lim_{x\to\pm\infty}\frac{\textsf{BIG}}{\textsf{small}}=$ \fillin{$\pm\infty$ (DNE, no HA)} \quad bigger degree on top\par
2) $\displaystyle\lim_{x\to\pm\infty}\frac{\textsf{small}}{\textsf{BIG}}=$ \fillin{$0$ (HA $y=0$)} \quad bigger degree on the bottom\par
3) $\displaystyle\lim_{x\to\pm\infty}\frac{\textsf{same}}{\textsf{same}}=$ \fillin{ratio of the leading coefficients} \quad same degree
\end{nrow}
\begin{nrow}{Examples}
1) $\displaystyle\lim_{x\to\infty}\frac{4x^3-x}{5-2x^2}=$ \fillin{DNE ($-\infty$)}\qquad
2) $\displaystyle\lim_{x\to-\infty}\frac{3x}{x^2+7}=$ \fillin{$0$}\par\smallskip
3) $\displaystyle\lim_{x\to\infty}\frac{6x^2+1}{3-2x^2}=$ \fillin{$\dfrac{6}{-2}=-3$}
\begin{workspace}[1.6cm]
1) degree $3>2$: grows without bound; $\frac{4x^3}{-2x^2}=-2x\to-\infty$. \;2) degree $1<2$: $0$. \;3) same degree $2$: $\frac{6}{-2}=-3$; HA $y=-3$.
\end{workspace}
\end{nrow}
\begin{nrow}{Summary}
\ifkey{\color{apnavy}Factor first: a cancelled factor gives a hole, a zero left in the denominator gives a vertical asymptote. For $x\to\pm\infty$, compare degrees: top bigger $\to\pm\infty$, bottom bigger $\to0$, equal $\to$ ratio of leading coefficients (the HA).}\else\vspace{2.4cm}\fi
\end{nrow}

%%DOC 1-4 | Continuity at a Point
\begin{center}\textsf{\small\color{apgrey}Objective: identify the types of discontinuities and prove continuity with the continuity test.}\end{center}
\begin{nrow}{What does it mean to be continuous?}
A graph is continuous if you can trace it \fillin{without lifting your pencil}.\par\smallskip
Are these graphs continuous or discontinuous?\par
\begin{tabular}{@{}cccc@{}}
\tinyg{xmin=-2.5,xmax=2.5,ymin=-0.5,ymax=3.5,width=3.6cm,height=3.3cm}{\addplot[apblue,very thick,domain=-2.4:2.4,samples=40]{x^2/2};} &
\tinyg{xmin=-2.5,xmax=3.5,ymin=-2.5,ymax=4.5,width=3.6cm,height=3.3cm}{\addplot[apred,very thick,domain=-2.4:1]{x};\addplot[apred,very thick,domain=1:3.4]{x+2};\op{1,1}\pt{1,3}} &
\tinyg{xmin=-3.5,xmax=3.5,ymin=-2.5,ymax=2.5,width=3.6cm,height=3.3cm}{\addplot[apgreen,very thick,domain=-3.4:3.4,samples=60]{1.5*sin(deg(x))};} &
\tinyg{xmin=-2.5,xmax=4,ymin=-4,ymax=4,width=3.6cm,height=3.3cm}{\addplot[apgold!80!black,very thick,domain=-2.4:0.75,samples=40]{1/(x-1)};\addplot[apgold!80!black,very thick,domain=1.25:3.9,samples=40]{1/(x-1)};\va{1}}\\
\fillin{continuous} & \fillin{discontinuous} & \fillin{continuous} & \fillin{discontinuous}
\end{tabular}
\end{nrow}
\begin{nrow}{Types of discontinuities}
\begin{tabular}{@{}ccc@{}}
\tinyg{xmin=-2.5,xmax=3.5,ymin=-1.5,ymax=4.5,width=4.6cm,height=3.9cm}{\addplot[apblue,very thick,domain=-2.4:3.4]{x+1};\op{1,2}\pt{1,-0.5}} &
\tinyg{xmin=-2.5,xmax=3.5,ymin=-1.5,ymax=4.5,width=4.6cm,height=3.9cm}{\addplot[apblue,very thick,domain=-2.4:1]{0.5*x+1};\addplot[apblue,very thick,domain=1:3.4]{0.5*x+2.5};\pt{1,1.5}\op{1,3}} &
\tinyg{xmin=-2.5,xmax=4,ymin=-1,ymax=5,width=4.6cm,height=3.9cm}{\addplot[apblue,very thick,domain=-2.4:0.55,samples=40]{1/(x-1)^2};\addplot[apblue,very thick,domain=1.45:3.9,samples=40]{1/(x-1)^2};\va{1}}\\
\fillin{Removable (hole)} & \fillin{Jump} & \fillin{Infinite (VA)}\\[2pt]
\footnotesize the limit \fillin{exists} but & \footnotesize the one-sided limits & \footnotesize the function goes\\
\footnotesize is not equal to $f(c)$ & \footnotesize are \fillin{different} & \footnotesize to \fillin{$\pm\infty$}
\end{tabular}\par\smallskip
An \textbf{endpoint} (like the start of $y=\sqrt x$ at $x=0$) only needs the one-sided limit from the \fillin{inside} of the domain.
\end{nrow}
\begin{nrow}{Continuity test}
$f$ is continuous at $x=c$ when
\[\fillin{$\displaystyle\lim_{x\to c^-}f(x)$}\;=\;\fillin{$\displaystyle\lim_{x\to c^+}f(x)$}\;=\;\fillin{$f(c)$}\]
\centering{\small\color{apgrey}\qquad left \hspace{2.4cm} right \hspace{2.2cm} actual height}
\end{nrow}
\begin{nrow}{Example 1}
$h(x)=\begin{cases}2x-1, & x\le3\\ x^2-4, & x>3\end{cases}$\qquad Prove that $h$ is continuous at $x=3$.
\begin{workspace}[2.6cm]
Left: $\lim_{x\to3^-}h(x)=2(3)-1=5$. \;Right: $\lim_{x\to3^+}h(x)=3^2-4=5$. \;Height: $h(3)=2(3)-1=5$.\\
$5=5=5$, so $h$ is continuous at $x=3$ because it passes the continuity test.
\end{workspace}
\end{nrow}
\begin{nrow}{Example 2}
For what value of $k$ is $g$ continuous at $x=1$? \qquad $g(x)=\begin{cases}3x^2-k, & x\le1\\ kx+5, & x>1\end{cases}$
\begin{workspace}[2.8cm]
Set left $=$ right at $x=1$: $3(1)^2-k=k(1)+5$, so $3-k=k+5$, $-2=2k$, $k=-1$.\\
Check: left $3-(-1)=4$, right $-1+5=4$, $g(1)=4$. Continuous.
\end{workspace}
\end{nrow}
\begin{nrow}{Summary}
\ifkey{\color{apnavy}Continuous means no hole, no jump and no asymptote at the point. To prove it, show left limit $=$ right limit $=f(c)$. To find a constant, set the two pieces equal at the suspect point and solve.}\else\vspace{2.4cm}\fi
\end{nrow}

%%DOC 1-5 | IVT and EVT
\goals{\item State the Intermediate Value Theorem and the Extreme Value Theorem, with their hypotheses.
\item Use the IVT to justify that an equation has a solution.
\item Use the EVT to justify that absolute extrema exist, and find them from graphs and tables.}
\vocab{Intermediate Value Theorem (IVT), Extreme Value Theorem (EVT), closed interval, absolute maximum, absolute minimum, hypothesis, conclusion.}

\subsection*{1. Intermediate Value Theorem}
\begin{keyidea}[IVT]
If $f$ is \fillin{continuous} on the \fillin{closed} interval $[a,b]$ and $k$ is any number \fillin{between $f(a)$ and $f(b)$}, then there is at least one $c$ in $[a,b]$ with \fillin{$f(c)=k$}.
\end{keyidea}
\gr{xmin=-0.5,xmax=4.5,ymin=-5.5,ymax=7.5,ytick distance=2,width=6cm,height=5cm}{\addplot[apblue,very thick,domain=0:4]{(x-1)*(x-2)*(x-3)+1};\ha{1}\pt{0,-5}\pt{4,7}}
\ex{Example 1}
Show that $x^3+2x=5$ has a solution in $[1,2]$.
\begin{workspace}[2.6cm]
$f(x)=x^3+2x$ is a polynomial, so it is continuous on $[1,2]$. $f(1)=3$, $f(2)=12$. Since $3<5<12$, by the IVT there is a $c$ in $[1,2]$ with $f(c)=5$.
\end{workspace}
\ex{Example 2 \textbullet\ Table}
$W$ is continuous with $W(0)=8$, $W(2)=3$, $W(5)=6$, $W(9)=1$. What is the least number of solutions of $W(t)=4$ on $[0,9]$?
\begin{workspace}[2.2cm]
$4$ is between $8$ and $3$, between $3$ and $6$, and between $6$ and $1$: at least three.
\end{workspace}

\subsection*{2. Extreme Value Theorem}
\begin{keyidea}[EVT]
If $f$ is \fillin{continuous} on the closed interval $[a,b]$, then $f$ has both an \fillin{absolute maximum} and an \fillin{absolute minimum} value on $[a,b]$. They occur at \fillin{endpoints} or at points where the graph turns (critical points).
\end{keyidea}

\ex{Example 3}
Find the absolute maximum and minimum of $f(x)=x^2-4x+5$ on $[0,3]$. Why must they exist?
\begin{workspace}[2.8cm]
$f$ is a polynomial, continuous on $[0,3]$, so by the EVT both exist. Candidates: vertex $x=2$, $f(2)=1$; endpoints $f(0)=5$, $f(3)=2$. Absolute maximum $5$ at $x=0$; absolute minimum $1$ at $x=2$.
\end{workspace}

\ex{Example 4 \textbullet\ When the theorems do not apply}
Explain why the conclusion can fail: (a) $g(x)=\frac1x$ on $[-1,1]$, IVT with $k=0$; (b) $h(x)=\frac1x$ on $(0,1]$, EVT.
\begin{workspace}[2.6cm]
(a) $g(-1)=-1$, $g(1)=1$, yet $g$ is never $0$: $g$ is not continuous on $[-1,1]$.\\
(b) $h$ has no absolute maximum (it grows without bound near $0$): the interval is not closed.
\end{workspace}

\begin{aptip}
An IVT or EVT justification names the hypothesis (``$f$ is continuous on $[a,b]$ because\dots''), the values used, and the conclusion.
\end{aptip}

\begin{quickcheck}
1. $f$ continuous, $f(1)=4$, $f(3)=4$. Does the IVT guarantee $f(c)=0$ on $[1,3]$? \fillin{No}.\\
2. Does a continuous function on $(0,5)$ always have an absolute maximum? \fillin{No; the interval must be closed}.
\end{quickcheck}

%%DOC 1-6 | Average Rate of Change
\goals{\item Compute an average rate of change from a formula, a table or a graph.
\item Interpret an average rate of change in context with units, including average velocity.
\item Use the difference quotient $\frac{f(a+h)-f(a)}{h}$.}
\vocab{average rate of change, secant line, difference quotient, average velocity, units.}

\begin{keyidea}
The average rate of change of $f$ on $[a,b]$ is \fillin{$\dfrac{f(b)-f(a)}{b-a}$}, the slope of the \fillin{secant line} through $(a,f(a))$ and $(b,f(b))$. On $[a,a+h]$ it is the \textbf{difference quotient} \fillin{$\dfrac{f(a+h)-f(a)}{h}$}.
\end{keyidea}

\ex{Example 1 \textbullet\ Formula}
Find the average rate of change of $f(x)=x^3-2x$ on $[-1,2]$, and simplify $\dfrac{f(1+h)-f(1)}{h}$ for $f(x)=x^2$.
\begin{workspace}[2.4cm]
$\frac{4-1}{3}=1$. \quad $\frac{(1+h)^2-1}{h}=\frac{2h+h^2}{h}=2+h$.
\end{workspace}

\ex{Example 2 \textbullet\ Table, with units}
\begin{center}\begin{tabular}{c|ccccc}$t$ (min)&$0$&$5$&$10$&$15$&$20$\\\hline$A(t)$ (L)&$40$&$52$&$61$&$66$&$68$\end{tabular}\end{center}
Find and interpret the average rate of change of $A$ on $[5,15]$.
\begin{workspace}[2.2cm]
$\frac{66-52}{10}=1.4$. From $t=5$ to $t=15$ minutes, the amount of water increased on average by $1.4$ liters per minute.
\end{workspace}

\ex{Example 3 \textbullet\ Average velocity from a graph}
\gr{xmin=-0.5,xmax=6.5,ymin=-0.5,ymax=8.5,ytick distance=2,width=6cm,height=5cm,xlabel={$t$},ylabel={$s$}}{\addplot[apblue,very thick] coordinates {(0,0)(2,6)(4,6)(6,2)};}
The position $s(t)$ (meters) of a particle is graphed. Find its average velocity on $[0,2]$, $[2,4]$ and $[0,6]$.
\begin{workspace}[2.2cm]
$\frac{6-0}{2}=3$ m/s. \quad $0$ m/s. \quad $\frac{2-0}{6}=\frac13$ m/s.
\end{workspace}

\begin{mistakes}
\begin{itemize}
\item Giving the change $f(b)-f(a)$ instead of the rate.
\item Forgetting units, or interpreting an average rate as a rate at an instant.
\end{itemize}
\end{mistakes}

\begin{quickcheck}
1. The average rate of change of $g(x)=\frac{12}{x}$ on $[2,6]$ is \fillin{$-1$}.\\
2. For $f(x)=x^2+x$, $\frac{f(2+h)-f(2)}{h}=$ \fillin{$5+h$}.
\end{quickcheck}

%%DOC 1-7 | Instantaneous Rates of Change
\goals{\item Define the derivative at a point as the limit of a difference quotient.
\item Compute derivatives of simple functions from the definition and estimate them from tables and graphs.
\item Write the equation of a tangent line, and recognize where a function is not differentiable.}
\vocab{instantaneous rate of change, derivative, $f'(a)$, tangent line, differentiable, corner, cusp, vertical tangent.}

\subsection*{1. From average to instantaneous}
\begin{keyidea}[Definition of the derivative]
\[f'(a)=\lim_{h\to0}\fillin{$\dfrac{f(a+h)-f(a)}{h}$}\qquad\text{or}\qquad f'(a)=\lim_{x\to a}\fillin{$\dfrac{f(x)-f(a)}{x-a}$}.\]
$f'(a)$ is the \fillin{instantaneous rate of change} of $f$ at $a$ and the \fillin{slope of the tangent line} at $(a,f(a))$.
\end{keyidea}
\gr{xmin=-0.5,xmax=3.5,ymin=-1,ymax=9.5,ytick distance=2,width=6cm,height=5.4cm}{\addplot[apblue,very thick,domain=0:3]{x^2};\addplot[apred!50,thick,domain=0:3.2]{4*x-3};\addplot[apred!75,thick,domain=0:3.2]{3*x-2};\addplot[apred,very thick,dashed,domain=0.2:3.2]{2*x-1};\pt{1,1}}
\begin{center}\small Secant lines from $(1,1)$ on $y=x^2$ approach the tangent line of slope $2$.\end{center}

\ex{Example 1 \textbullet\ From the definition}
Find $f'(x)$ for $f(x)=x^2$, $g(x)=x^3$ and $k(x)=\frac1x$.
\begin{workspace}[4.4cm]
$\frac{(x+h)^2-x^2}{h}=2x+h\to2x$: $f'(x)=2x$.\\
$\frac{(x+h)^3-x^3}{h}=3x^2+3xh+h^2\to3x^2$: $g'(x)=3x^2$.\\
$\frac{\frac1{x+h}-\frac1x}{h}=\frac{-1}{x(x+h)}\to-\frac1{x^2}$: $k'(x)=-\frac1{x^2}$.
\end{workspace}

\ex{Example 2 \textbullet\ Tangent line}
Write the equation of the tangent line to $y=x^3$ at $x=2$.
\begin{workspace}[2cm]
Point $(2,8)$; slope $g'(2)=3(4)=12$. \ $y-8=12(x-2)$.
\end{workspace}

\ex{Example 3 \textbullet\ Estimating from a table}
\begin{center}\begin{tabular}{c|ccccc}$t$ (s)&$0$&$2$&$4$&$6$&$8$\\\hline$s(t)$ (m)&$0$&$9$&$20$&$35$&$52$\end{tabular}\end{center}
Estimate $s'(4)$ and interpret it.
\begin{workspace}[2.2cm]
Use the closest values on both sides: $\frac{s(6)-s(2)}{6-2}=\frac{35-9}{4}=6.5$. At $t=4$ s the object moves at about $6.5$ m/s.
\end{workspace}

\subsection*{2. Differentiability}
\begin{keyidea}
If $f$ is differentiable at $a$, then $f$ is \fillin{continuous} at $a$. The converse is \fillin{false}.\\
$f$ is \textbf{not} differentiable where its graph has a \fillin{corner}, a \fillin{cusp}, a \fillin{vertical tangent}, or a \fillin{discontinuity}.
\end{keyidea}
\begin{center}
\begin{tikzpicture}
\begin{axis}[gridaxis,xmin=-1.6,xmax=1.6,ymin=-0.5,ymax=1.8,width=4.2cm,height=3.8cm,title={\scriptsize corner: $|x|$},xtick=\empty,ytick=\empty]
\addplot[apblue,very thick,domain=-1.5:1.5]{abs(x)};\end{axis}\end{tikzpicture}
\begin{tikzpicture}
\begin{axis}[gridaxis,xmin=-1.6,xmax=1.6,ymin=-0.5,ymax=1.8,width=4.2cm,height=3.8cm,title={\scriptsize cusp: $x^{2/3}$},xtick=\empty,ytick=\empty]
\addplot[apblue,very thick,domain=-1.5:1.5,samples=200]{abs(x)^(2/3)};\end{axis}\end{tikzpicture}
\begin{tikzpicture}
\begin{axis}[gridaxis,xmin=-1.6,xmax=1.6,ymin=-1.4,ymax=1.4,width=4.2cm,height=3.8cm,title={\scriptsize vertical tangent: $\sqrt[3]{x}$},xtick=\empty,ytick=\empty]
\addplot[apblue,very thick,domain=0:1.5,samples=200]{x^(1/3)};\addplot[apblue,very thick,domain=-1.5:0,samples=200]{-(-x)^(1/3)};\end{axis}\end{tikzpicture}
\end{center}

\ex{Example 4 \textbullet\ A piecewise function}
Is $f(x)=\begin{cases}x^2, & x\le1\\ 2x-1, & x>1\end{cases}$ differentiable at $x=1$?
\begin{workspace}[2.6cm]
Continuous: both pieces give $1$ at $x=1$. Slopes: the left piece $x^2$ has slope $2(1)=2$; the right piece has slope $2$. The slopes match, so $f$ is differentiable at $x=1$ with $f'(1)=2$.
\end{workspace}

\begin{quickcheck}
1. $\lim_{h\to0}\frac{(3+h)^2-9}{h}$ is $f'(3)$ for $f(x)=$ \fillin{$x^2$}, and equals \fillin{$6$}.\\
2. $f(x)=|x-2|$ is continuous at $2$ but \fillin{not differentiable} there.
\end{quickcheck}

%%DOC 1-8 | MVT
\goals{\item State the Mean Value Theorem and check its hypotheses.
\item Find the value(s) of $c$ guaranteed by the MVT.
\item Apply the MVT to tables and contexts.}
\vocab{Mean Value Theorem (MVT), differentiable, open interval, secant slope, tangent slope.}

\begin{keyidea}[Mean Value Theorem]
If $f$ is \fillin{continuous} on $[a,b]$ and \fillin{differentiable} on $(a,b)$, then there is at least one $c$ in $(a,b)$ such that
\[f'(c)=\fillin{$\dfrac{f(b)-f(a)}{b-a}$}.\]
Some tangent line is \fillin{parallel} to the secant line: at some instant, the instantaneous rate of change equals the average rate of change.
\end{keyidea}
\gr{xmin=-0.5,xmax=3.5,ymin=-0.5,ymax=9.5,ytick distance=2,width=6cm,height=5.2cm}{\addplot[apblue,very thick,domain=0.5:3.2]{x^2};\addplot[apred,thick,domain=0.6:3.3]{4*x-3};\addplot[apgreen,thick,dashed,domain=0.8:3.2]{4*x-4};\pt{1,1}\pt{3,9}\pt{2,4}}
\begin{center}\small $f(x)=x^2$ on $[1,3]$: secant slope $4$; the tangent at $c=2$ has slope $4$ too.\end{center}

\ex{Example 1 \textbullet\ Finding $c$}
Find the value(s) of $c$ guaranteed by the MVT for (a) $f(x)=x^2$ on $[1,3]$; (b) $g(x)=x^3$ on $[0,3]$; (c) $k(x)=\frac1x$ on $[1,4]$. (Use $f'(x)=2x$, $g'(x)=3x^2$, $k'(x)=-\frac1{x^2}$ from Lesson 1-7.)
\begin{workspace}[3.6cm]
(a) $\frac{9-1}{2}=4=2c$: $c=2$.\\
(b) $\frac{27-0}{3}=9=3c^2$: $c=\sqrt3\approx1.732$ (the negative root is not in $(0,3)$).\\
(c) $\frac{\frac14-1}{3}=-\frac14=-\frac1{c^2}$: $c=2$.
\end{workspace}

\ex{Example 2 \textbullet\ When the MVT does not apply}
Explain why the MVT does not apply: (a) $f(x)=|x|$ on $[-1,1]$; (b) $g(x)=\frac1x$ on $[-1,1]$.
\begin{workspace}[2.4cm]
(a) Not differentiable at $x=0$ (corner). The secant slope is $0$, yet no tangent is horizontal.\\
(b) Not continuous on $[-1,1]$ (undefined at $0$).
\end{workspace}

\ex{Example 3 \textbullet\ A table in context}
A car's position $s(t)$ (miles) is differentiable, with $s(0)=0$ and $s(2)=130$ ($t$ in hours). Must the car's speed have been exactly $65$ mph at some time?
\begin{workspace}[2.4cm]
$s$ is differentiable, so continuous; the MVT applies on $[0,2]$: $s'(c)=\frac{130-0}{2}=65$ for some $c$ in $(0,2)$. Yes: at some instant the velocity was $65$ mph.
\end{workspace}

\begin{mistakes}
\begin{itemize}
\item Forgetting to check differentiability (corners, cusps) or continuity.
\item Giving a $c$ that is not inside the open interval.
\end{itemize}
\end{mistakes}

\begin{quickcheck}
1. For $f(x)=x^2$ on $[0,4]$, the MVT gives $c=$ \fillin{$2$}.\\
2. If $f(2)=f(5)$ and $f$ satisfies the hypotheses, then $f'(c)=$ \fillin{$0$} for some $c$ in $(2,5)$.
\end{quickcheck}

%%DOC 1-9 | IVT, EVT, and MVT
\goals{\item Choose the right existence theorem for a question.
\item Check the hypotheses of the IVT, EVT and MVT from formulas, tables and graphs.
\item Write complete AP-style justifications.}
\vocab{existence theorem, hypothesis, conclusion, justification.}

\subsection*{1. Which theorem?}
\begin{keyidea}
\renewcommand{\arraystretch}{1.4}
\begin{tabular}{@{}p{1.3cm}p{5cm}p{6.3cm}@{}}
& \textbf{Hypothesis} & \textbf{Conclusion}\\\hline
IVT & \fillin{continuous on $[a,b]$} & $f(c)=k$ for every $k$ between \fillin{$f(a)$ and $f(b)$}\\
EVT & \fillin{continuous on $[a,b]$} & $f$ has an absolute \fillin{maximum and minimum} on $[a,b]$\\
MVT & continuous on $[a,b]$ and \fillin{differentiable on $(a,b)$} & $f'(c)=$ \fillin{$\frac{f(b)-f(a)}{b-a}$} for some $c$ in $(a,b)$
\end{tabular}
\end{keyidea}
Questions about a \fillin{value of $f$} point to the IVT; about the \fillin{largest or smallest value}, the EVT; about a \fillin{rate (value of $f'$)}, the MVT.

\ex{Example 1 \textbullet\ A table (AP style)}
A differentiable function $f$ has the values below.
\begin{center}\begin{tabular}{c|ccccc}$x$&$0$&$2$&$5$&$7$&$10$\\\hline$f(x)$&$4$&$-2$&$3$&$9$&$5$\end{tabular}\end{center}
(a) Must $f(c)=0$ for some $c$ in $[0,10]$? (b) Must $f'(c)=3$ for some $c$ in $(5,7)$? (c) Must $f$ have an absolute maximum on $[0,10]$? (d) Must $f'(c)=0$ for some $c$ in $(0,10)$?
\begin{workspace}[5.2cm]
$f$ is differentiable, hence continuous everywhere.\\
(a) Yes: on $[0,2]$, $f(0)=4>0>-2=f(2)$; by the IVT there is $c$ with $f(c)=0$ (also on $[2,5]$).\\
(b) Yes: $\frac{f(7)-f(5)}{7-5}=\frac{6}{2}=3$; by the MVT, $f'(c)=3$ for some $c$ in $(5,7)$.\\
(c) Yes: $f$ is continuous on the closed interval $[0,10]$; by the EVT it has an absolute maximum.\\
(d) Yes. $f(0)=4$, and since $3<4<9$, the IVT gives a $b$ in $(5,7)$ with $f(b)=4$. On $[0,b]$ the MVT gives $f'(c)=\frac{f(b)-f(0)}{b-0}=\frac{4-4}{b}=0$ for some $c$ in $(0,b)$.
\end{workspace}

\ex{Example 2 \textbullet\ Hypotheses fail}
\gr{xmin=-0.5,xmax=4.5,ymin=-0.5,ymax=4.5,width=5.5cm,height=4.6cm}{\addplot[apblue,very thick,domain=0:2]{x+1};\addplot[apblue,very thick,domain=2:4]{3-0.5*(x-2)^2};\pt{0,1}\pt{4,1}\pt{2,3}}
The graph of $f$ on $[0,4]$ is shown (a corner at $x=2$). Which theorems apply on $[0,4]$?
\begin{workspace}[2.6cm]
$f$ is continuous on $[0,4]$: the IVT and EVT apply. $f$ is not differentiable at $x=2$ (corner), so the MVT does not apply on $[0,4]$ (it does apply on $[2,4]$).
\end{workspace}

\begin{aptip}
Never apply a theorem without stating why its hypotheses hold. On tables, ``$f$ is differentiable, so it is continuous'' covers both continuity and differentiability.
\end{aptip}

\begin{quickcheck}
1. ``Is there a time when the temperature was exactly $70^\circ$?'' uses the \fillin{IVT}.\\
2. ``Is there a time when the temperature was rising at $2^\circ$ per hour?'' uses the \fillin{MVT}.
\end{quickcheck}
